\documentclass[12pt,a4paper]{article}

\usepackage[T2A]{fontenc}
\usepackage[utf8]{inputenc}
\usepackage[russian]{babel}
\usepackage{geometry}
\usepackage{setspace}
\usepackage{indentfirst}
\usepackage{hyperref}
\usepackage{longtable}
\usepackage{array}

\geometry{top=2cm,bottom=1.5cm,left=3cm,right=2cm}
\onehalfspacing

\hypersetup{
  colorlinks=true,
  linkcolor=black,
  urlcolor=blue,
  pdftitle={Пояснительная записка: Эмулятор CPU 6502},
  pdfauthor={Ступин Александр Юрьевич}
}

\begin{document}

\begin{titlepage}
\begin{center}
\small
\textbf{Муниципальное бюджетное общеобразовательное учреждение\\
«Гимназия №2 «Квантор»}}\\[0.7cm]

\textbf{Индивидуальный проект}\\
по предмету: \textit{Информатика}\\
тип проекта: \textit{учебный, практико-ориентированный}\\[1.2cm]

{\LARGE \textbf{Разработка учебного веб-приложения\\
«Эмулятор CPU 6502 с ассемблером, консолью и пошаговой трассировкой»}}\\[1.5cm]

\begin{flushright}
\begin{tabular}{l}
Выполнил: Ступин Александр Юрьевич \\
Класс: 10 «В» \\
Руководитель проекта: Пименова Анна Николаевна
\end{tabular}
\end{flushright}

\vfill
\textbf{Коломна --- 2026}
\end{center}
\end{titlepage}

\tableofcontents
\newpage

\section{Введение}

\subsection{Актуальность темы}
В школьном курсе информатики учащиеся в основном работают с языками высокого уровня. При этом понимание того, как реально исполняются инструкции процессора, как меняются регистры и память, часто остаётся фрагментарным. Тема архитектуры ЭВМ изучается теоретически, без постоянной практической визуализации.

Актуальность проекта заключается в создании интерактивного учебного инструмента, который позволяет в реальном времени наблюдать работу условной модели CPU 6502: загрузку программы, выполнение команд, изменение регистров, флагов и памяти по шагам. Такой формат делает сложную тему более понятной и практически значимой.

\subsection{Проблема}
Существующие эмуляторы процессоров либо избыточно сложны для учебных целей школы, либо не имеют удобного интерфейса, ориентированного на пошаговый анализ. В результате ученику трудно связать теорию (инструкции, флаги, переходы) с практикой исполнения программы.

\subsection{Гипотеза}
Если разработать учебное веб-приложение, объединяющее эмулятор CPU 6502, ассемблер и пошаговую визуализацию состояния процессора, то понимание учащимися принципов низкоуровневого программирования и архитектуры процессора станет более глубоким и осознанным.

\subsection{Объект и предмет исследования}
\textbf{Объект исследования} --- процесс обучения архитектуре ЭВМ и основам ассемблерного программирования в школьном курсе информатики.

\textbf{Предмет исследования} --- методы разработки и педагогическая эффективность интерактивного программного продукта для визуализации исполнения инструкций CPU.

\subsection{Цель проекта}
Разработать и апробировать учебное веб-приложение «Эмулятор CPU 6502», позволяющее вводить программу на ассемблере, запускать её и пошагово анализировать исполнение команд по состоянию регистров, флагов и памяти.

\subsection{Задачи проекта}
\begin{enumerate}
  \item Изучить архитектуру и историю процессора 6502, определить учебно значимые элементы его модели.
  \item Изучить инструменты разработки: Python, WebSocket, React, TypeScript.
  \item Спроектировать архитектуру приложения (frontend + backend + протокол обмена).
  \item Реализовать эмулятор CPU и ассемблер на Python.
  \item Реализовать WebSocket-сервер для взаимодействия клиентской и серверной части.
  \item Реализовать пользовательский интерфейс с консолью, справкой по инструкциям и step-view трассы.
  \item Провести тестирование, оценить результаты и сформулировать направления дальнейшего развития.
\end{enumerate}

\subsection{Проектный продукт}
Проектный продукт --- учебное веб-приложение «Эмулятор CPU 6502», включающее:
\begin{itemize}
  \item текстовый редактор ассемблерного кода;
  \item модуль сборки/запуска программы;
  \item консоль ввода и вывода значений;
  \item пошаговый просмотр исполнения (STEP) с навигацией стрелками;
  \item отображение регистров A, X, Y, PC, OP и блока MEMORY;
  \item светлую и тёмную темы интерфейса;
  \item справочный раздел по инструкциям.
\end{itemize}

\section{Основная часть}

\subsection{Краткая историческая справка о CPU 6502}
Микропроцессор MOS Technology 6502 был представлен в 1975 году и стал одним из самых влиятельных процессоров эпохи ранних персональных компьютеров. Он использовался, в частности, в Apple II, Commodore и ряде игровых систем. Простота архитектуры и низкая стоимость обеспечили массовое распространение и сделали 6502 удобной учебной платформой.

Для образовательных задач 6502 важен как пример классической архитектуры с ограниченным, но выразительным набором регистров и флагов, где легко проследить причинно-следственную связь между командой и изменением состояния системы.

\subsection{Теоретические основы, используемые в проекте}
В проекте использованы следующие понятия:
\begin{itemize}
  \item архитектура процессора: регистры, флаги, память, счётчик команд;
  \item ассемблерный язык и машинные инструкции;
  \item условные переходы и относительная адресация;
  \item клиент-серверное взаимодействие в реальном времени (WebSocket);
  \item компонентный подход к разработке интерфейса.
\end{itemize}

\subsection{Концепция проекта и целевая аудитория}
Веб-приложение «CPU6502\_ts» проектируется как современная цифровая платформа для интерактивной эмуляции классического 8-битного процессора MOS 6502. Концептуально проект объединяет два уровня:
\begin{itemize}
  \item точную низкоуровневую модель исполнения инструкций процессора;
  \item удобный браузерный интерфейс уровня современного developer-tool.
\end{itemize}

Целевая аудитория:
\begin{itemize}
  \item учащиеся и студенты технических направлений;
  \item энтузиасты ретро-компьютинга;
  \item начинающие и практикующие разработчики, изучающие low-level программирование.
\end{itemize}

Платформа ориентирована на пользователей, которым важны быстрый отклик интерфейса, наглядность состояния CPU и возможность проводить эксперименты с ассемблерным кодом без установки сложного локального ПО.

\subsection{Выбор технологий и full-stack реализация}
\textbf{Python} выбран для backend-части из-за читаемости, удобства тестирования и скорости реализации алгоритмов эмуляции.\newline
\textbf{WebSocket} обеспечивает двунаправленный канал в реальном времени между интерфейсом и эмулятором.\newline
\textbf{React + TypeScript} формируют типобезопасный и масштабируемый интерфейс с компонентной архитектурой.\newline
\textbf{Vite} используется для быстрой сборки и разработки frontend.\newline
\textbf{Caddy} применяется в продакшен-развёртывании как reverse-proxy с автоматической выдачей HTTPS-сертификатов.

Дополнительно применяются:
\begin{itemize}
  \item JSON-протокол сообщений между frontend и backend;
  \item переменные окружения для конфигурации и определения WebSocket-endpoint;
  \item Git-репозиторий для контроля версий и публикации обновлений.
\end{itemize}

\subsection{Архитектура программного решения}
Система реализована как распределённое приложение из двух частей:
\begin{enumerate}
  \item \textbf{Backend} (Python): CPU-модель, ассемблер, выполнение программы, формирование трассы шагов, WebSocket API.
  \item \textbf{Frontend} (React + TypeScript): редактор кода, консоль, отображение байткода и итогового состояния, step-view по шагам исполнения.
\end{enumerate}

Передача данных организована в формате JSON:
\begin{itemize}
  \item \texttt{assemble} --- сборка исходного текста в байты;
  \item \texttt{run} --- выполнение программы с передачей входной очереди;
  \item \texttt{result} --- возврат трассы, памяти, выходных значений и финального состояния.
\end{itemize}

\subsection{Пользовательский интерфейс и UX-принципы}
Интерфейс построен по принципу «всё важное на одном экране». Пользователь одновременно видит:
\begin{itemize}
  \item редактор исходного кода;
  \item панель запуска и управления;
  \item консоль ввода/вывода;
  \item итоговое состояние процессора;
  \item пошаговую визуализацию исполнения.
\end{itemize}

Принципы UX:
\begin{itemize}
  \item минимализм и техническая ясность без перегрузки декоративными элементами;
  \item логическая иерархия блоков (сначала код и запуск, затем результаты и трасса);
  \item высокая читаемость данных в тёмной теме;
  \item адаптивность для разных экранов (ПК/ноутбук/планшет);
  \item мгновенная обратная связь при действиях пользователя.
\end{itemize}

\subsection{Описание функциональных модулей}
\subsubsection{Модуль эмулятора CPU}
Эмулятор реализует:
\begin{itemize}
  \item регистры \texttt{A, X, Y, SP, PC, P};
  \item флаги состояния \texttt{C, Z, N, V} и служебные;
  \item память 64 КБ;
  \item цикл исполнения инструкций с учётом циклов.
\end{itemize}

\subsubsection{Модуль ассемблера}
Ассемблер выполняет двухпроходный разбор: сначала определяет адреса меток, затем формирует машинный код. Поддерживаются инструкции загрузки, арифметики, логики, работы с памятью, переходов и служебные команды.

\subsubsection{Модуль пошаговой визуализации}
В интерфейсе реализован режим Step, где для каждого шага доступны:
\begin{itemize}
  \item номер шага \texttt{STEP i/n};
  \item значения регистров \texttt{A, X, Y, PC, OP};
  \item состояние выполнения (running/halted/error);
  \item блок \texttt{MEMORY} с занятыми/изменёнными ячейками.
\end{itemize}

\subsubsection{Консольный модуль взаимодействия}
Отдельный блок консоли обеспечивает учебный сценарий, близкий к реальной работе с программой:
\begin{itemize}
  \item добавление входных значений в очередь (hex);
  \item просмотр информационных сообщений запуска;
  \item отображение выходных значений программы;
  \item очистка очереди и повторный запуск с новыми параметрами.
\end{itemize}

\subsection{Методы исследования и разработки}
Применялись методы:
\begin{itemize}
  \item анализ предметной области и литературных источников;
  \item модульное проектирование;
  \item итеративная разработка (планирование --- реализация --- тестирование --- улучшение интерфейса);
  \item функциональное тестирование сценариев выполнения программ.
\end{itemize}

\subsection{Этапы выполнения проекта}
\begin{enumerate}
  \item Формулировка темы, целей и требований к продукту.
  \item Изучение архитектуры 6502 и необходимого набора инструкций.
  \item Реализация Python-модулей CPU и ассемблера.
  \item Реализация WebSocket-сервера.
  \item Создание frontend-интерфейса React + TypeScript.
  \item Интеграция и отладка обмена клиент-сервер.
  \item UX-доработки: темы, консоль, step-view, справка.
  \item Подготовка проекта к защите.
\end{enumerate}

\subsection{Технико-экономическое обоснование}
Для данного учебного программного проекта отдельная смета материального производства не требуется, так как продукт цифровой и реализуется на общедоступном программном стеке.

Используемые ресурсы:
\begin{itemize}
  \item персональный компьютер разработчика;
  \item бесплатные/учебные инструменты разработки (Python, Node.js, Git);
  \item хостинг/VPS для публикации веб-приложения.
\end{itemize}

Вывод: проект имеет низкую стоимость внедрения и может быть тиражирован в образовательной среде без значительных затрат.

\subsection{Охрана труда, техника безопасности и информационная безопасность}
Так как проект программный, ключевые меры связаны с безопасной организацией работы за ПК:
\begin{itemize}
  \item соблюдение санитарно-гигиенических норм работы с экраном;
  \item регулярные перерывы и контроль зрительной нагрузки;
  \item корректная посадка и организация рабочего места.
\end{itemize}

В части информационной безопасности:
\begin{itemize}
  \item ограничение публичных входных точек сервера;
  \item контроль формата входных данных;
  \item безопасное хранение ключей и паролей (вне исходного кода);
  \item регулярное обновление программных зависимостей.
\end{itemize}

\subsection{Результаты проекта и их оценка}
В результате получено рабочее приложение, которое:
\begin{itemize}
  \item собирает и запускает пользовательские программы;
  \item показывает пошаговое выполнение;
  \item предоставляет визуальный контроль состояния регистров и памяти;
  \item поддерживает учебный сценарий объяснения переходов и флагов.
\end{itemize}

По итогам тестирования достигнуты заявленные цели и решены поставленные задачи. Продукт применим как в рамках учебного занятия, так и для самостоятельной подготовки.

\subsection{Образовательная и практическая значимость}
Реализованный эмулятор выполняет одновременно учебную и демонстрационно-портфельную функции:
\begin{itemize}
  \item повышает вовлечённость в изучение архитектуры ЭВМ за счёт интерактивности;
  \item позволяет наглядно связать теорию (регистры, флаги, переходы) и практику;
  \item формирует у автора компетенции full-stack разработки;
  \item может использоваться как цифровая витрина технического проекта.
\end{itemize}

\subsection{Перспективы развития}
Планируемые направления развития платформы:
\begin{itemize}
  \item добавление breakpoints и расширенного пошагового отладчика;
  \item визуализация памяти в hex-представлении;
  \item поддержка прерываний и cycle-accurate режима эмуляции;
  \item эмуляция периферийных устройств (видео/звук);
  \item сохранение программ в localStorage и/или базе данных;
  \item библиотека готовых примеров (демо, учебные кейсы, мини-игры);
  \item развитие клиентской версии эмулятора на TypeScript.
\end{itemize}

\section{Заключение}
Цель проекта --- создание учебного веб-приложения «Эмулятор CPU 6502» --- достигнута. Разработаны backend и frontend части, реализованы механизмы ассемблирования, выполнения и пошагового анализа программ. Полученный результат подтверждает выдвинутую гипотезу: интерактивный формат повышает наглядность и качество понимания архитектуры процессора.

Итоговый продукт представляет собой технологически целостную full-stack платформу, сочетающую исторически значимую архитектуру 6502 и современные веб-практики. Проект имеет практическую значимость для школьного курса информатики, может быть использован как учебный инструмент и как элемент инженерного портфолио автора.

\newpage
\section*{Список использованных источников}
\addcontentsline{toc}{section}{Список использованных источников}

\begin{enumerate}
  \item Федеральный государственный образовательный стандарт среднего общего образования (ФГОС СОО).
  \item Python Software Foundation. Python 3 Documentation. URL: \url{https://docs.python.org/3/} (дата обращения: 01.03.2026).
  \item Ramalho L. \textit{Fluent Python}. 2nd ed. O'Reilly Media, 2022.
  \item Downey A. \textit{Think Python}. 2nd ed. O'Reilly Media, 2015.
  \item Beazley D., Jones B. \textit{Python Cookbook}. 3rd ed. O'Reilly Media, 2013.
  \item React Team. React Documentation. URL: \url{https://react.dev/} (дата обращения: 01.03.2026).
  \item Microsoft. TypeScript Documentation. URL: \url{https://www.typescriptlang.org/docs/} (дата обращения: 01.03.2026).
  \item Vite Team. Vite Guide. URL: \url{https://vite.dev/guide/} (дата обращения: 01.03.2026).
  \item Banks A., Porcello E. \textit{Learning React}. 3rd ed. O'Reilly Media, 2024.
  \item Cherny B. \textit{Programming TypeScript}. O'Reilly Media, 2019.
  \item Leventhal L. \textit{6502 Assembly Language Programming}. Osborne/McGraw-Hill, 1979.
  \item Zaks R. \textit{Programming the 6502}. Sybex, 1981.
  \item 8bitworkshop. 6502 and retro development resources. URL: \url{https://8bitworkshop.com/} (дата обращения: 01.03.2026).
  \item Masswerk. 6502 Instruction Set Reference. URL: \url{https://www.masswerk.at/6502/} (дата обращения: 01.03.2026).
  \item MOS Technology. \textit{MCS6500 Microcomputer Family Programming Manual}. 1976.
  \item MDN Web Docs. WebSocket API. URL: \url{https://developer.mozilla.org/docs/Web/API/WebSocket} (дата обращения: 01.03.2026).
  \item IETF. RFC 6455: The WebSocket Protocol. URL: \url{https://datatracker.ietf.org/doc/html/rfc6455} (дата обращения: 01.03.2026).
  \item Caddy Server Documentation. URL: \url{https://caddyserver.com/docs/} (дата обращения: 01.03.2026).
  \item Pro Git Book. URL: \url{https://git-scm.com/book/} (дата обращения: 01.03.2026).
\end{enumerate}

\newpage
\section*{План защиты проекта (6--8 минут)}
\addcontentsline{toc}{section}{План защиты проекта (6--8 минут)}

\begin{enumerate}
  \item Вступление: тема и актуальность (до 1 мин).
  \item Проблема, цель, задачи, гипотеза (1 мин).
  \item Архитектура решения и выбор технологий (1.5 мин).
  \item Демонстрация продукта: запуск, console, step-view (2 мин).
  \item Результаты, практическая значимость, перспективы (1 мин).
  \item Ответы на вопросы комиссии (оставшееся время).
\end{enumerate}

\newpage
\section*{Текст защиты проекта}
\addcontentsline{toc}{section}{Текст защиты проекта}

Здравствуйте, уважаемая комиссия.

Меня зовут Ступин Александр Юрьевич, я ученик 10 «В» класса МБОУ «Гимназия №2 «Квантор» города Коломна. Представляю индивидуальный проект по информатике на тему: «Разработка учебного веб-приложения “Эмулятор CPU 6502 с ассемблером, консолью и пошаговой трассировкой”».

Актуальность проекта в том, что при изучении программирования мы часто видим только код высокого уровня и не понимаем, что происходит внутри процессора. Из-за этого тема архитектуры ЭВМ остаётся сложной и абстрактной.

Проблема проекта --- отсутствие простого учебного инструмента, который наглядно показывает исполнение команд, изменение регистров, флагов и памяти.

Цель проекта --- создать интерактивное веб-приложение для изучения принципов работы CPU на примере 6502.

Для достижения цели я решил задачи: изучил архитектуру 6502, реализовал эмулятор и ассемблер на Python, связал backend и frontend через WebSocket, и сделал интерфейс на React + TypeScript с консолью и пошаговым режимом.

Архитектура решения: пользователь вводит код в редакторе, фронтенд отправляет его на Python-сервер, сервер собирает и исполняет программу, а затем возвращает трассу шагов и итоговое состояние. В интерфейсе можно переключать шаги стрелками и смотреть значения A, X, Y, PC, OP и блок MEMORY.

Практический результат --- готовое учебное приложение, которое можно использовать на уроках информатики и для самостоятельной подготовки. Оно помогает понять, как именно инструкции влияют на состояние процессора.

Гипотеза подтвердилась: интерактивный формат действительно делает тему архитектуры ЭВМ более понятной и наглядной.

Перспективы: расширение набора инструкций, добавление задач с автопроверкой и улучшение визуализации сложных механизмов, например стека и прерываний.

Спасибо за внимание. Готов ответить на вопросы.

\end{document}
